\documentclass[10pt,oneside,openany]{memoir}

\iffalse
%font style 1
\usepackage[p,osf]{scholax}
\usepackage{amsmath,amsthm,amssymb}
\usepackage[scaled=1.075,ncf,vvarbb]{newtxmath}
\fi

%font style 2
\usepackage[T1]{fontenc}
\usepackage{amsmath}
\usepackage[fixamsmath]{mathtools}  % Extension to amsmath
\usepackage{amsthm}
\usepackage{amssymb}
\renewcommand*{\mathbf}[1]{\varmathbb{#1}}
\usepackage{newpxtext}
\usepackage{eulerpx}
\usepackage{mathrsfs}

\usepackage{eucal}
\usepackage{datetime}
    \newdateformat{specialdate}{\THEYEAR\ \monthname\ \THEDAY}
\usepackage[margin=1.5in]{geometry}
\usepackage{fancyhdr}
    \pagestyle{fancy}
    \fancyhf{}
    \cfoot{\scriptsize\thepage}
    \fancyhead[R]{\scalebox{0.8}{\nouppercase{\rightmark}}}
    \fancyhead[L]{\scalebox{0.8}{\nouppercase{\leftmark}}}

    \fancypagestyle{plain}{%
        \fancyhf{}
        \cfoot{\scriptsize\thepage}
        \fancyhead[R]{\scalebox{0.8}{\phantom{a}}}
        \fancyhead[L]{\scalebox{0.8}{Applied Analysis}}
    }
\usepackage{keytheorems}
    \newkeytheoremstyle{def}{
    inherit-style=definition,
    spaceabove=1.5em,
    spacebelow=1.5em,
    postheadspace=0.5em,
    }
    \newkeytheoremstyle{thm}{
    spaceabove=1.5em,
    spacebelow=1.5em,
    postheadspace=0.5em,
    }
    \newkeytheorem{theorem}[
    name = Theorem,
    style=thm,
    numberwithin = section,
    ]
    \newkeytheorem{theorem*}[
    name=Theorem,
    style=thm,
    numbered=false
    ]
    \newkeytheorem{proposition}[
    name = Proposition,
    style=thm,
    numberwithin = section,
    sibling=theorem
    ]
    \newkeytheorem{lemma}[
    name = Lemma,
    style=thm,
    numberwithin = section,
    sibling=theorem
    ]
    \newkeytheorem{corollary}[
    name = Corollary,
    style=thm,
    numberwithin = section,
    sibling=theorem
    ]
    \newkeytheorem{definition}[
    name = Definition,
    style=def,
    numberwithin = section,
    style = definition
    ]
    \newkeytheorem{example}[
    name = Example,
    style=def,
    numberwithin = section,
    style = definition
    ]
    \newkeytheorem{remark}[
    name=Remark,
    style=remark, 
    numbered=false  
    ]
    \newkeytheorem{question}[
    name=Question,
    style=remark, 
    numbered=false  
    ]
    \newkeytheorem{exercise}[
    name = Exercise,
    style=thm
    ]
    \newkeytheorem{axiom}[name = Axiom]
\usepackage{enumitem}
\usepackage[explicit]{titlesec}
    \titleformat{\chapter}[display]
    {\bfseries\LARGE\raggedright}
    {Chapter {\thechapter}}
    {.1ex minus .1ex}
    {\Large #1}
    \titlespacing{\chapter}
    {3pc}{*3}{40pt}[3pc]

    \titleformat{\section}[block]
    {\normalfont\bfseries\large}
    {}{0pt}
    {\fbox{\strut \S\ \thesection.\ #1}}
    \titlespacing{\section}
    {0pt}{3ex plus .1ex minus .2ex}{3ex plus .1ex minus .2ex}
    \setlength{\fboxsep}{4pt}   % padding inside the box
    \setlength{\fboxrule}{0.6pt} % border thickness
    
    \titleformat{name=\section,numberless}[block]
    {\normalfont\bfseries\large\raggedright}
    {}{0pt}
    {\fbox{\strut #1}}
\usepackage[utf8x]{inputenc}
\usepackage{tikz}
\usepackage{tikz-cd}
\usepackage{float}
\usetikzlibrary{patterns}
\usetikzlibrary{positioning,arrows.meta}
\usetikzlibrary{decorations.pathreplacing}
\usepackage{wasysym}
\usepackage{hyperref}
\usepackage{ulem}
\usepackage{pgfplots}
\usepackage[svgnames]{xcolor}
\usetikzlibrary{patterns}
\usepgfplotslibrary{fillbetween}
\pgfplotsset{compat=1.18}

\renewcommand{\headrulewidth}{0pt}
\renewcommand{\footrulewidth}{0pt}
%%%%%%%%%%%%%%%%%%%%%%%%%%%%%%%%%%%%%%%%%%%%%%%%%%%%%%%%%%%%%
%%%%%%%%%%%%%%%%%%%%%%%%%%%%%%%%%%%%%%%%%%%%%%%%%%%%%%%%%%%%%
% ============================================================
% makros.tex — reorganized, full restored version
% ============================================================

% ============================================================
% Sha / Cyrillic support notes
% ============================================================

%to make the correct symbol for Sha
%\newcommand\cyr{%
%\renewcommand\rmdefault{wncyr}%
%\renewcommand\sfdefault{wncyss}%
%\renewcommand\encodingdefault{OT2}%
%\normalfont \selectfont} \DeclareTextFontCommand{\textcyr}{\cyr}

% ============================================================
% Math operators
% ============================================================

\DeclareMathOperator{\ab}{ab}
\DeclareMathOperator{\ad}{ad}
\DeclareMathOperator{\adj}{adj}
\DeclareMathOperator{\alg}{alg}
\DeclareMathOperator{\Alt}{Alt}
\DeclareMathOperator{\Ann}{Ann}
\DeclareMathOperator{\arith}{arith}
\DeclareMathOperator{\Aut}{Aut}
\DeclareMathOperator{\Be}{B}
\DeclareMathOperator{\Bd}{Bd}
\DeclareMathOperator{\card}{card}
\DeclareMathOperator{\Char}{char}
\DeclareMathOperator{\csp}{csp}
\DeclareMathOperator{\codim}{codim}
\DeclareMathOperator{\coker}{coker}
\DeclareMathOperator{\coh}{H}
\DeclareMathOperator{\compl}{compl}
\DeclareMathOperator{\conj}{conj}
\DeclareMathOperator{\cont}{cont}
\DeclareMathOperator{\Cov}{Cov}
\DeclareMathOperator{\crys}{crys}
\DeclareMathOperator{\Crys}{Crys}
\DeclareMathOperator{\cusp}{cusp}
\DeclareMathOperator{\diag}{diag}
\DeclareMathOperator{\diam}{diam}
\DeclareMathOperator{\Dom}{Dom}
\DeclareMathOperator{\disc}{disc}
\DeclareMathOperator{\dist}{dist}
\DeclareMathOperator{\Div}{div}
\DeclareMathOperator{\dR}{dR}
\DeclareMathOperator{\Eis}{Eis}
\DeclareMathOperator{\End}{End}
\DeclareMathOperator{\ev}{ev}
\DeclareMathOperator{\eval}{eval}
\DeclareMathOperator{\Eq}{Eq}
\DeclareMathOperator{\Ext}{Ext}
\DeclareMathOperator{\Fil}{Fil}
\DeclareMathOperator{\Fitt}{Fitt}
\DeclareMathOperator{\Frob}{Frob}
\DeclareMathOperator{\G}{G}
\DeclareMathOperator{\Gal}{Gal}
\DeclareMathOperator{\GL}{GL}
\DeclareMathOperator{\Gr}{Gr}
\DeclareMathOperator{\Graph}{Graph}
\DeclareMathOperator{\GSp}{GSp}
\DeclareMathOperator{\GUn}{GU}
\DeclareMathOperator{\Hom}{Hom}
\DeclareMathOperator{\id}{id}
\DeclareMathOperator{\Id}{Id}
\DeclareMathOperator{\Ik}{Ik}
\DeclareMathOperator{\IM}{Im}
\DeclareMathOperator{\Image}{im}
\DeclareMathOperator{\Ind}{Ind}
\DeclareMathOperator{\Inf}{inf}
\DeclareMathOperator{\ins}{ins}
\DeclareMathOperator{\inv}{inv}
\DeclareMathOperator{\Isom}{Isom}
\DeclareMathOperator{\J}{J}
\DeclareMathOperator{\Jac}{Jac}
\DeclareMathOperator{\lcm}{lcm}
\DeclareMathOperator{\length}{length}
\DeclareMathOperator*{\limit}{limit}
\DeclareMathOperator{\Log}{Log}
\DeclareMathOperator{\M}{M}
\DeclareMathOperator{\Mat}{Mat}
\DeclareMathOperator{\N}{N}
\DeclareMathOperator{\Nm}{Nm}
\DeclareMathOperator{\NIk}{N-Ik}
\DeclareMathOperator{\NSK}{N-SK}
\DeclareMathOperator{\new}{new}
\DeclareMathOperator{\obj}{obj}
\DeclareMathOperator{\old}{old}
\DeclareMathOperator{\ord}{ord}
\DeclareMathOperator{\Or}{O}
\DeclareMathOperator{\op}{op}
\DeclareMathOperator{\out}{out}
\DeclareMathOperator{\PGL}{PGL}
\DeclareMathOperator{\PGSp}{PGSp}
\DeclareMathOperator{\rank}{rank}
\DeclareMathOperator{\Ran}{Ran}
\DeclareMathOperator{\rel}{Rel}
\DeclareMathOperator{\Real}{Re}
\DeclareMathOperator{\RES}{res}
\DeclareMathOperator{\Res}{Res}
%\DeclareMathOperator{\Sha}{\textcyr{Sh}}
\DeclareMathOperator{\Sel}{Sel}
\DeclareMathOperator{\semi}{ss}
\DeclareMathOperator{\sgn}{sign}
\DeclareMathOperator{\SK}{SK}
\DeclareMathOperator{\SL}{SL}
\DeclareMathOperator{\SO}{SO}
\DeclareMathOperator{\Sp}{Sp}
\DeclareMathOperator{\Span}{span}
\DeclareMathOperator{\Spec}{Spec}
\DeclareMathOperator{\spin}{spin}
\DeclareMathOperator{\st}{st}
\DeclareMathOperator{\St}{St}
\DeclareMathOperator{\SUn}{SU}
\DeclareMathOperator{\supp}{supp}
\DeclareMathOperator{\Sup}{sup}
\DeclareMathOperator{\Sym}{Sym}
\DeclareMathOperator{\Tam}{Tam}
\DeclareMathOperator{\tors}{tors}
\DeclareMathOperator{\tr}{tr}
\DeclareMathOperator{\Tr}{Tr}
\DeclareMathOperator{\un}{un}
\DeclareMathOperator{\Un}{U}
\DeclareMathOperator{\val}{val}
\DeclareMathOperator{\vol}{vol}

% ============================================================
% General aliases and short macros
% ============================================================

\newcommand{\absgal}{\G_{\bbQ}}
\newcommand{\Loc}{L_{\operatorname{loc}}}
\renewcommand{\k}{\kappa}
\newcommand{\Ff}{F_{f}}
%\newcommand{\ts}{\,^{t}\!}
\newcommand{\lat}{\mathscr{L}}
\newcommand{\ov}[1]{\overline{#1}}
\newcommand{\up}{\upsilon}
\newcommand{\e}{\epsilon}
\newcommand{\tac}{\textasteriskcentered}

%mahesh macros
\newcommand{\tm}{\textrm}

\newcommand{\bone}{\mathbf{1}}
\newcommand{\sign}{\mathrm{sign}}
\newcommand{\eps}{\varepsilon}
\newcommand{\textui}[1]{\underline{\textit{#1}}}

%\newcommand{\ov}{\overline}
%\newcommand{\un}{\underline}
\newcommand{\fin}{\mathrm{fin}}
\newcommand{\nl}{\newline \mbox{}}
\newcommand{\h}[1]{\hspace{#1pt}}
\newcommand{\mtext}[1]{\hspace{6pt}\text{#1}\hspace{6pt}}
\newcommand{\lnorm}{\left\lVert}
\newcommand{\rnorm}{\right\rVert}
\newcommand{\ds}{\displaystyle}
\newcommand{\ts}{\textstyle}
\newcommand{\sfrac}[2]{{}^{#1}\mskip -5mu/\mskip -3mu_{#2}}

% ============================================================
% Category-style operators and contoured variants
% ============================================================

\DeclareMathOperator{\Sets}{S \mkern1.04mu e \mkern1.04mu t \mkern1.04mu s}
    \newcommand{\cSets}{\scalebox{1.02}{\contour{black}{$\Sets$}}}
    
\DeclareMathOperator{\Groups}{G \mkern1.04mu r \mkern1.04mu o \mkern1.04mu u \mkern1.04mu p \mkern1.04mu s}
    \newcommand{\cGroups}{\scalebox{1.02}{\contour{black}{$\Groups$}}}

\DeclareMathOperator{\TTop}{T \mkern1.04mu o \mkern1.04mu p}
    \newcommand{\cTop}{\scalebox{1.02}{\contour{black}{$\TTop$}}}

\DeclareMathOperator{\Htp}{H \mkern1.04mu t \mkern1.04mu p}
    \newcommand{\cHtp}{\scalebox{1.02}{\contour{black}{$\Htp$}}}

\DeclareMathOperator{\Mod}{M \mkern1.04mu o \mkern1.04mu d}
    \newcommand{\cMod}{\scalebox{1.02}{\contour{black}{$\Mod$}}}

\DeclareMathOperator{\Ab}{A \mkern1.04mu b}
    \newcommand{\cAb}{\scalebox{1.02}{\contour{black}{$\Ab$}}}

\DeclareMathOperator{\Rings}{R \mkern1.04mu i \mkern1.04mu n \mkern1.04mu g \mkern1.04mu s}
    \newcommand{\cRings}{\scalebox{1.02}{\contour{black}{$\Rings$}}}

\DeclareMathOperator{\ComRings}{C \mkern1.04mu o \mkern1.04mu m \mkern1.04mu R \mkern1.04mu i \mkern1.04mu n \mkern1.04mu g \mkern1.04mu s}
    \newcommand{\cComRings}{\scalebox{1.05}{\contour{black}{$\ComRings$}}}

\DeclareMathOperator{\hHom}{H \mkern1.04mu o \mkern1.04mu m}
    \newcommand{\cHom}{\scalebox{1.02}{\contour{black}{$\hHom$}}}

% ============================================================
% Mathcal
% ============================================================

\newcommand{\cA}{\mathcal{A}}
\newcommand{\cB}{\mathcal{B}}
\newcommand{\cC}{\mathcal{C}}
\newcommand{\cD}{\mathcal{D}}
\newcommand{\cE}{\mathcal{E}}
\newcommand{\cF}{\mathcal{F}}
\newcommand{\cG}{\mathcal{G}}
\newcommand{\cH}{\mathcal{H}}
\newcommand{\cI}{\mathcal{I}}
\newcommand{\cJ}{\mathcal{J}}
\newcommand{\cK}{\mathcal{K}}
\newcommand{\cL}{\mathcal{L}}
\newcommand{\cM}{\mathcal{M}}
\newcommand{\cN}{\mathcal{N}}
\newcommand{\cO}{\mathcal{O}}
\newcommand{\cP}{\mathcal{P}}
\newcommand{\cQ}{\mathcal{Q}}
\newcommand{\cR}{\mathcal{R}}
\newcommand{\cS}{\mathcal{S}}
\newcommand{\cT}{\mathcal{T}}
\newcommand{\cU}{\mathcal{U}}
\newcommand{\cV}{\mathcal{V}}
\newcommand{\cW}{\mathcal{W}}
\newcommand{\cX}{\mathcal{X}}
\newcommand{\cY}{\mathcal{Y}}
\newcommand{\cZ}{\mathcal{Z}}

% ============================================================
% Mathscrl% 
%============================================================

\newcommand{\scra}{\mathscr{a}}
\newcommand{\scrA}{\mathscr{A}}
\newcommand{\scrb}{\mathscr{b}}
\newcommand{\scrB}{\mathscr{B}}
\newcommand{\scrc}{\mathscr{c}}
\newcommand{\scrC}{\mathscr{C}}
\newcommand{\scrd}{\mathscr{d}}
\newcommand{\scrD}{\mathscr{D}}
\newcommand{\scre}{\mathscr{e}}
\newcommand{\scrE}{\mathscr{E}}
\newcommand{\scrf}{\mathscr{f}}
\newcommand{\scrF}{\mathscr{F}}
\newcommand{\scrg}{\mathscr{g}}
\newcommand{\scrG}{\mathscr{G}}
\newcommand{\scrh}{\mathscr{h}}
\newcommand{\scrH}{\mathscr{H}}
\newcommand{\scri}{\mathscr{i}}
\newcommand{\scrI}{\mathscr{I}}
\newcommand{\scrj}{\mathscr{j}}
\newcommand{\scrJ}{\mathscr{J}}
\newcommand{\scrk}{\mathscr{k}}
\newcommand{\scrK}{\mathscr{K}}
\newcommand{\scrl}{\mathscr{l}}
\newcommand{\scrL}{\mathscr{L}}
\newcommand{\scrm}{\mathscr{m}}
\newcommand{\scrM}{\mathscr{M}}
\newcommand{\scrn}{\mathscr{n}}
\newcommand{\scrN}{\mathscr{N}}
\newcommand{\scro}{\mathscr{o}}
\newcommand{\scrO}{\mathscr{O}}
\newcommand{\scrp}{\mathscr{p}}
\newcommand{\scrP}{\mathscr{P}}
\newcommand{\scrq}{\mathscr{q}}
\newcommand{\scrQ}{\mathscr{Q}}
\newcommand{\scrr}{\mathscr{r}}
\newcommand{\scrR}{\mathscr{R}}
\newcommand{\scrs}{\mathscr{s}}
\newcommand{\scrS}{\mathscr{S}}
\newcommand{\scrt}{\mathscr{t}}
\newcommand{\scrT}{\mathscr{T}}
\newcommand{\scru}{\mathscr{u}}
\newcommand{\scrU}{\mathscr{U}}
\newcommand{\scrv}{\mathscr{v}}
\newcommand{\scrV}{\mathscr{V}}
\newcommand{\scrw}{\mathscr{w}}
\newcommand{\scrW}{\mathscr{W}}
\newcommand{\scrx}{\mathscr{x}}
\newcommand{\scrX}{\mathscr{X}}
\newcommand{\scry}{\mathscr{y}}
\newcommand{\scrY}{\mathscr{Y}}
\newcommand{\scrz}{\mathscr{z}}
\newcommand{\scrZ}{\mathscr{Z}}

% ============================================================
% Mathfrak (missing \fi)
% ============================================================

\newcommand{\fa}{\mathfrak{a}}
\newcommand{\fA}{\mathfrak{A}}
\newcommand{\fb}{\mathfrak{b}}
\newcommand{\fB}{\mathfrak{B}}
\newcommand{\fc}{\mathfrak{c}}
\newcommand{\fC}{\mathfrak{C}}
\newcommand{\fd}{\mathfrak{d}}
\newcommand{\fD}{\mathfrak{D}}
\newcommand{\fe}{\mathfrak{e}}
\newcommand{\fE}{\mathfrak{E}}
\newcommand{\ff}{\mathfrak{f}}
\newcommand{\fF}{\mathfrak{F}}
\newcommand{\fg}{\mathfrak{g}}
\newcommand{\fG}{\mathfrak{G}}
\newcommand{\fh}{\mathfrak{h}}
\newcommand{\fH}{\mathfrak{H}}
\newcommand{\fI}{\mathfrak{I}}
\newcommand{\fj}{\mathfrak{j}}
\newcommand{\fJ}{\mathfrak{J}}
\newcommand{\fk}{\mathfrak{k}}
\newcommand{\fK}{\mathfrak{K}}
\newcommand{\fl}{\mathfrak{l}}
\newcommand{\fL}{\mathfrak{L}}
\newcommand{\fm}{\mathfrak{m}}
\newcommand{\fM}{\mathfrak{M}}
\newcommand{\fn}{\mathfrak{n}}
\newcommand{\fN}{\mathfrak{N}}
\newcommand{\fo}{\mathfrak{o}}
\newcommand{\fO}{\mathfrak{O}}
\newcommand{\fp}{\mathfrak{p}}
\newcommand{\fP}{\mathfrak{P}}
\newcommand{\fq}{\mathfrak{q}}
\newcommand{\fQ}{\mathfrak{Q}}
\newcommand{\fr}{\mathfrak{r}}
\newcommand{\fR}{\mathfrak{R}}
\newcommand{\fs}{\mathfrak{s}}
\newcommand{\fS}{\mathfrak{S}}
\newcommand{\ft}{\mathfrak{t}}
\newcommand{\fT}{\mathfrak{T}}
\newcommand{\fu}{\mathfrak{u}}
\newcommand{\fU}{\mathfrak{U}}
\newcommand{\fv}{\mathfrak{v}}
\newcommand{\fV}{\mathfrak{V}}
\newcommand{\fw}{\mathfrak{w}}
\newcommand{\fW}{\mathfrak{W}}
\newcommand{\fx}{\mathfrak{x}}
\newcommand{\fX}{\mathfrak{X}}
\newcommand{\fy}{\mathfrak{y}}
\newcommand{\fY}{\mathfrak{Y}}
\newcommand{\fz}{\mathfrak{z}}
\newcommand{\fZ}{\mathfrak{Z}}

% ============================================================
% Mathbf
% ============================================================

\newcommand{\bfA}{\mathbf{A}}
\newcommand{\bfB}{\mathbf{B}}
\newcommand{\bfC}{\mathbf{C}}
\newcommand{\bfD}{\mathbf{D}}
\newcommand{\bfE}{\mathbf{E}}
\newcommand{\bfF}{\mathbf{F}}
\newcommand{\bfG}{\mathbf{G}}
\newcommand{\bfH}{\mathbf{H}}
\newcommand{\bfI}{\mathbf{I}}
\newcommand{\bfJ}{\mathbf{J}}
\newcommand{\bfK}{\mathbf{K}}
\newcommand{\bfL}{\mathbf{L}}
\newcommand{\bfM}{\mathbf{M}}
\newcommand{\bfN}{\mathbf{N}}
\newcommand{\bfO}{\mathbf{O}}
\newcommand{\bfP}{\mathbf{P}}
\newcommand{\bfQ}{\mathbf{Q}}
\newcommand{\bfR}{\mathbf{R}}
\newcommand{\bfS}{\mathbf{S}}
\newcommand{\bfT}{\mathbf{T}}
\newcommand{\bfU}{\mathbf{U}}
\newcommand{\bfV}{\mathbf{V}}
\newcommand{\bfW}{\mathbf{W}}
\newcommand{\bfX}{\mathbf{X}}
\newcommand{\bfY}{\mathbf{Y}}
\newcommand{\bfZ}{\mathbf{Z}}

\newcommand{\bfa}{\mathbf{a}}
\newcommand{\bfb}{\mathbf{b}}
\newcommand{\bfc}{\mathbf{c}}
\newcommand{\bfd}{\mathbf{d}}
\newcommand{\bfe}{\mathbf{e}}
\newcommand{\bff}{\mathbf{f}}
\newcommand{\bfg}{\mathbf{g}}
\newcommand{\bfh}{\mathbf{h}}
\newcommand{\bfi}{\mathbf{i}}
\newcommand{\bfj}{\mathbf{j}}
\newcommand{\bfk}{\mathbf{k}}
\newcommand{\bfl}{\mathbf{l}}
\newcommand{\bfm}{\mathbf{m}}
\newcommand{\bfn}{\mathbf{n}}
\newcommand{\bfo}{\mathbf{o}}
\newcommand{\bfp}{\mathbf{p}}
\newcommand{\bfq}{\mathbf{q}}
\newcommand{\bfr}{\mathbf{r}}
\newcommand{\bfs}{\mathbf{s}}
\newcommand{\bft}{\mathbf{t}}
\newcommand{\bfu}{\mathbf{u}}
\newcommand{\bfv}{\mathbf{v}}
\newcommand{\bfw}{\mathbf{w}}
\newcommand{\bfx}{\mathbf{x}}
\newcommand{\bfy}{\mathbf{y}}
\newcommand{\bfz}{\mathbf{z}}

% ============================================================
% Blackboard bold
% ============================================================

\newcommand{\bbA}{\mathbb{A}}
\newcommand{\bbB}{\mathbb{B}}
\newcommand{\bbC}{\mathbb{C}}
\newcommand{\bbD}{\mathbb{D}}
\newcommand{\bbE}{\mathbb{E}}
\newcommand{\bbF}{\mathbb{F}}
\newcommand{\bbG}{\mathbb{G}}
\newcommand{\bbH}{\mathbb{H}}
\newcommand{\bbI}{\mathbb{I}}
\newcommand{\bbJ}{\mathbb{J}}
\newcommand{\bbK}{\mathbb{K}}
\newcommand{\bbL}{\mathbb{L}}
\newcommand{\bbM}{\mathbb{M}}
\newcommand{\bbN}{\mathbb{N}}
\newcommand{\bbO}{\mathbb{O}}
\newcommand{\bbP}{\mathbb{P}}
\newcommand{\bbQ}{\mathbb{Q}}
\newcommand{\bbR}{\mathbb{R}}
\newcommand{\bbS}{\mathbb{S}}
\newcommand{\bbT}{\mathbb{T}}
\newcommand{\bbU}{\mathbb{U}}
\newcommand{\bbV}{\mathbb{V}}
\newcommand{\bbW}{\mathbb{W}}
\newcommand{\bbX}{\mathbb{X}}
\newcommand{\bbY}{\mathbb{Y}}
\newcommand{\bbZ}{\mathbb{Z}}

\newcommand{\Bmu}{\mbox{$\raisebox{-0.59ex}
  {$l$}\hspace{-0.18em}\mu\hspace{-0.88em}\raisebox{-0.98ex}{\scalebox{2}
  {$\color{white}.$}}\hspace{-0.416em}\raisebox{+0.88ex}
  {$\color{white}.$}\hspace{0.46em}$}{}}  %need graphicx and xcolor. this produces blackboard bold mu 

% ============================================================
% Greek-letter aliases
% ============================================================

\newcommand{\jota}{\jmath}

% ============================================================
% Matrices
% ============================================================

\newcommand{\bmat}{\left( \begin{matrix}}
\newcommand{\emat}{\end{matrix} \right)}

\newcommand{\bbmat}{\left[ \begin{matrix}}
\newcommand{\ebmat}{\end{matrix} \right]}

\newcommand{\pmat}{\left( \begin{smallmatrix}}
\newcommand{\epmat}{\end{smallmatrix} \right)}

\newcommand{\mat}[4]{\begin{pmatrix}{#1}&{#2}\\{#3}&{#4}\end{pmatrix}}

% ============================================================
% Residues and averaged integrals
% ============================================================

\newcommand{\res}[1]{\underset{#1}{\RES}\,}

\makeatletter
\newcommand*{\mint}[1]{%
  % #1: overlay symbol
  \mint@l{#1}{}%
}
\newcommand*{\mint@l}[2]{%
  % #1: overlay symbol
  % #2: limits
  \@ifnextchar\limits{%
    \mint@l{#1}%
  }{%
    \@ifnextchar\nolimits{%
      \mint@l{#1}%
    }{%
      \@ifnextchar\displaylimits{%
        \mint@l{#1}%
      }{%
        \mint@s{#2}{#1}%
      }%
    }%
  }%
}
\newcommand*{\mint@s}[2]{%
  % #1: limits
  % #2: overlay symbol
  \@ifnextchar_{%
    \mint@sub{#1}{#2}%
  }{%
    \@ifnextchar^{%
      \mint@sup{#1}{#2}%
    }{%
      \mint@{#1}{#2}{}{}%
    }%
  }%
}
\def\mint@sub#1#2_#3{%
  \@ifnextchar^{%
    \mint@sub@sup{#1}{#2}{#3}%
  }{%
    \mint@{#1}{#2}{#3}{}%
  }%
}
\def\mint@sup#1#2^#3{%
  \@ifnextchar_{%
    \mint@sup@sub{#1}{#2}{#3}%
  }{%
    \mint@{#1}{#2}{}{#3}%
  }%
}
\def\mint@sub@sup#1#2#3^#4{%
  \mint@{#1}{#2}{#3}{#4}%
}
\def\mint@sup@sub#1#2#3_#4{%
  \mint@{#1}{#2}{#4}{#3}%
}
\newcommand*{\mint@}[4]{%
  % #1: \limits, \nolimits, \displaylimits
  % #2: overlay symbol: -, =, ...
  % #3: subscript
  % #4: superscript
  \mathop{}%
  \mkern-\thinmuskip
  \mathchoice{%
    \mint@@{#1}{#2}{#3}{#4}%
        \displaystyle\textstyle\scriptstyle
  }{%
    \mint@@{#1}{#2}{#3}{#4}%
        \textstyle\scriptstyle\scriptstyle
  }{%
    \mint@@{#1}{#2}{#3}{#4}%
        \scriptstyle\scriptscriptstyle\scriptscriptstyle
  }{%
    \mint@@{#1}{#2}{#3}{#4}%
        \scriptscriptstyle\scriptscriptstyle\scriptscriptstyle
  }%
  \mkern-\thinmuskip
  \int#1%
  \ifx\\#3\\\else_{#3}\fi
  \ifx\\#4\\\else^{#4}\fi  
}
\newcommand*{\mint@@}[7]{%
  % #1: limits
  % #2: overlay symbol
  % #3: subscript
  % #4: superscript
  % #5: math style
  % #6: math style for overlay symbol
  % #7: math style for subscript/superscript
  \begingroup
    \sbox0{$#5\int\m@th$}%
    \sbox2{$#5\int_{}\m@th$}%
    \dimen2=\wd0 %
    % => \dimen2 = width of \int
    \let\mint@limits=#1\relax
    \ifx\mint@limits\relax
      \sbox4{$#5\int_{\kern1sp}^{\kern1sp}\m@th$}%
      \ifdim\wd4>\wd2 %
        \let\mint@limits=\nolimits
      \else
        \let\mint@limits=\limits
      \fi
    \fi
    \ifx\mint@limits\displaylimits
      \ifx#5\displaystyle
        \let\mint@limits=\limits
      \fi
    \fi
    \ifx\mint@limits\limits
      \sbox0{$#7#3\m@th$}%
      \sbox2{$#7#4\m@th$}%
      \ifdim\wd0>\dimen2 %
        \dimen2=\wd0 %
      \fi
      \ifdim\wd2>\dimen2 %
        \dimen2=\wd2 %
      \fi
    \fi
    \rlap{%
      $#5%
        \vcenter{%
          \hbox to\dimen2{%
            \hss
            $#6{#2}\m@th$%
            \hss
          }%
        }%
      $%
    }%
  \endgroup
}
\newcommand{\avg}{\mint{-}}
\makeatother

% ============================================================
% Comments and drafting helpers
% ============================================================

%Comments
\newcommand{\com}[1]{\vspace{5 mm}\par \noindent
\marginpar{\textsc{Comment}} \framebox{\begin{minipage}[c]{0.95
\textwidth} \tt #1 \end{minipage}}\vspace{5 mm}\par}

% ============================================================
% Arrows
% ============================================================

\newcommand{\hooktwoheadrightarrow}{%
  \hookrightarrow\mathrel{\mspace{-15mu}}\rightarrow
}

\makeatletter
\newcommand{\xhooktwoheadrightarrow}[2][]{%
  \lhook\joinrel
  \ext@arrow 0359\rightarrowfill@ {#1}{#2}%
  \mathrel{\mspace{-15mu}}\rightarrow
}
\makeatother

\makeatletter
\newcommand*{\da@rightarrow}{\mathchar"0\hexnumber@\symAMSa 4B }
\newcommand*{\da@leftarrow}{\mathchar"0\hexnumber@\symAMSa 4C }
\newcommand*{\xdashrightarrow}[2][]{%
  \mathrel{%
    \mathpalette{\da@xarrow{#1}{#2}{}\da@rightarrow{\,}{}}{}%
  }%
}
\newcommand{\xdashleftarrow}[2][]{%
  \mathrel{%
    \mathpalette{\da@xarrow{#1}{#2}\da@leftarrow{}{}{\,}}{}%
  }%
}
\newcommand*{\da@xarrow}[7]{%
  % #1: below
  % #2: above
  % #3: arrow left
  % #4: arrow right
  % #5: space left 
  % #6: space right
  % #7: math style 
  \sbox0{$\ifx#7\scriptstyle\scriptscriptstyle\else\scriptstyle\fi#5#1#6\m@th$}%
  \sbox2{$\ifx#7\scriptstyle\scriptscriptstyle\else\scriptstyle\fi#5#2#6\m@th$}%
  \sbox4{$#7\dabar@\m@th$}%
  \dimen@=\wd0 %
  \ifdim\wd2 >\dimen@
    \dimen@=\wd2 %   
  \fi
  \count@=2 %
  \def\da@bars{\dabar@\dabar@}%
  \@whiledim\count@\wd4<\dimen@\do{%
    \advance\count@\@ne
    \expandafter\def\expandafter\da@bars\expandafter{%
      \da@bars
      \dabar@ 
    }%
  }%  
  \mathrel{#3}%
  \mathrel{%   
    \mathop{\da@bars}\limits
    \ifx\\#1\\%
    \else
      _{\copy0}%
    \fi
    \ifx\\#2\\%
    \else
      ^{\copy2}%
    \fi
  }%   
  \mathrel{#4}%
}
\makeatother

% ============================================================
% Relation and delimiter spacing
% ============================================================

\renewcommand{\geq}{\geqslant}
\renewcommand{\leq}{\leqslant}
\newcommand{\midd}{\hspace{4pt}\middle|\hspace{4pt}}

% ============================================================
% Left Right Updated Deliminators
% ============================================================

\makeatletter
\NewDocumentCommand\Left{O{}}{%
    \IfValueTF{#1}{\notblank{#1}{\mathopen\bBigg@{#1}}{\left}}{\left}}
\NewDocumentCommand\Right{O{}}{%
    \IfValueTF{#1}{\notblank{#1}{\mathopen\bBigg@{#1}}{\right}}{\right}}
\makeatother

% ============================================================
% Restrictions
% ============================================================

\newcommand{\littletaller}{\mathchoice{\vphantom{\big|}}{}{}{}}
\newcommand\restr[2]{{% we make the whole thing an ordinary symbol
  \left.\kern-\nulldelimiterspace % automatically resize the bar with \right
  #1 % the function
  \littletaller % pretend it's a little taller at normal size
  \right|_{#2} % this is the delimiter
  }}

% ============================================================
% Chapter title formatting
% ============================================================

\newcommand{\chnum}{\titleformat
{\chapter} % command
[display] % shape
{\centering} % format
{\Huge \color{black} \shadowbox{\thechapter}} % label
{-0.5em} % sep (space between the number and title)
{\LARGE \color{black} \underline} % before-code
}

\newcommand{\chunnum}{\titleformat
{\chapter} % command
[display] % shape
{} % format
{} % label
{0em} % sep
{ \begin{flushright} \begin{tabular}{r}  \Huge \color{black}
} % before-code
[
\end{tabular} \end{flushright} \normalsize
] % after-code
}



\begin{document}
\begin{center}
    {\large Math 540 }\\[.175in]
    {Name:} {\underline{Gianluca Crescenzo\hspace*{2in}}}\\[0.15in]
    \end{center}
    \vspace{4pt}
%%%%%%%%%%%%%%%%%%%%%%%%%%%%%%%%%%%%%%%%%%%%%%%%%%
%%%%%%%%%%%%%%%%%%%%%%%%%%%%%%%%%%%%%%%%%%%%%%%%%%
%%%%%%%%%%%%%%%%%%%%%%%%%%%%%%%%%%%%%%%%%%%%%%%%%%
%%%%%%%%%%%%%%%%%%%%%%%%%%%%%%%%%%%%%%%%%%%%%%%%%%
\section*{Week 7 Problems}
\begin{exercise}[manual-num={1}]
Suppose that $S,T$ are bounded linear operators on a Hilbert space $H$. Suppose $ST$ is compact. Must either $S$ or $T$ be compact? Either prove or give an explicit counter-example using $H = \ell^2$.
\end{exercise}
    \begin{proof}
        Define $S:\ell^2 \rightarrow \ell^2$ by:
            \begin{equation*}
            \begin{split}
                (x_1,x_2,x_3,x_4,...) \mapsto (x_1,0,x_2,0,...).
            \end{split}
            \end{equation*}
        Define $T:\ell^2 \rightarrow \ell^2$ by:
            \begin{equation*}
            \begin{split}
                (x_1,x_2,x_3,x_4,...) \mapsto (0,x_2,0,x_3,...).
            \end{split}
            \end{equation*}
        Neither $S$ nor $T$ are compact, but $S \circ T$ is compact as it is the zero operator.
    \end{proof}

\begin{exercise}[manual-num={4}]
Let $(\lambda_n)$ be a bounded sequence of complex numbers. Show that the bounded linear operator $T:\ell^2 \rightarrow \ell^2$ defined by:
    \begin{equation*}
    \begin{split}
        T e^n = \lambda_n e^n
    \end{split}
    \end{equation*}
(and extended linearly) is compact if $\lambda_n \rightarrow 0$.
\end{exercise}
    \begin{proof}
        Define $T_k:\ell^2 \rightarrow \ell^2$ by $T_kx := \sum_{n = 1}^k \lambda_n (x,e_n)e_n$. Then each $T_k$ is compact, and:
            \begin{equation*}
            \begin{split}
                \lnorm (T-T_k)(x) \rnorm_{\ell^2}^2 
                & = \lnorm \sum_{n > k}\lambda_n (x,e^n)e^n \rnorm_{\ell^2}^2 \\
                & \leq \sum_{n > k}|\lambda_n|^2 |(x,e^n)|^2 \lnorm e^n \rnorm^2 \\
                & = \sum_{n > k}|\lambda_n|^2 |(x,e^n)|^2 \\
                & \leq \sum_{n > k}\lnorm \lambda_n \rnorm_{\ell^\infty}^2 |(x,e^n)|^2 \\
                & = \lnorm \lambda_n \rnorm_{\ell^\infty}^2 \sum_{n > k}|(x,e^n)|^2 \\
                & \leq \lnorm \lambda_n \rnorm_{\ell^\infty}^2 \lnorm x \rnorm_{\ell^2}^2.
            \end{split}
            \end{equation*}
        Square rooting both sides, then taking the supremum over all $x \in \ell^2$, $\lnorm x \rnorm_{\ell^2} = 1$ gives $\lnorm T - T_k \rnorm \leq \lnorm \lambda_n \rnorm_{\ell^\infty}$. Since $\lambda_n \rightarrow 0$, it must be the case that $\lnorm T - T_k \rnorm \rightarrow 0$. Hence $T_k \rightarrow T$. Since the set of all compact operators is closed in $\cL(\ell^2)$, it must be the case that $T$ is compact.
    \end{proof}

\begin{exercise}[manual-num={7}]
Let $X$ denote an inner-product space over $\bfC$, and let $T:X \rightarrow X$ be a linear operator which is self-adjoint. Prove the following facts concerning the eigenvalues, eigenspaces, and invariant subspaces of $T$:
    \begin{enumerate}[label = (\alph*),itemsep=1pt,topsep=3pt]
        \item Show that all eigenvalues of $T$ are real.
        \item Show that any two eigenvectors of $T$ which correspond to distinct eigenvalues must be orthogonal.
        \item Suppose $Y \subset X$ is a $T$-invariant subspace of $X$. Show that $Y^\perp$ is also $T$-invariant.
    \end{enumerate}
\end{exercise}
    \begin{proof}
        (a) Suppose $Tx = \lambda x$ for some nonzero $x \in X$. Since $T$ is self-adjoint, we have $(Tx,x) = (x,Tx)$; i.e., $(\lambda x , x) = (x , \lambda x)$. This is equivalent to $\lambda (x,x) = \overline{\lambda}(x,x)$. Rearranging gives $(\lambda - \overline{\lambda})(x,x) = 0$. Since we assumed $x\neq 0$, it must be the case that $\lambda - \overline{\lambda} = 0$; i.e., $\lambda = \overline{\lambda}$. This can only be true if $\lambda \in \bfR$, hence any eigenvalue of $T$ must be real.

        (b) Let $\lambda_1,\lambda_2$ be distinct eigenvalues with corresponding eigenvectors $u_1,u_2$. Since $T$ is self-adjoint, we have $(Tu_1,u_2) = (u_1,Tu_2)$. This is equivalent to $(\lambda_1 u_1,u_2) = (u_1,\lambda_2 u_2)$. Thus $\lambda_1(u_1,u_2) = \lambda_2(u_1,u_2)$. Rearranging gives $(\lambda_1 - \lambda_2)(u_1,u_2) = 0$. Since $\lambda_1 \neq \lambda_2$, it must be the case that $(u_1,u_2) = 0$. Thus $u_1 \perp u_2$.

        (c) Let $z \in Y^\perp$. Then for any $y \in Y$ we have $(Tz,y) = (z,Ty)$. Since $Y$ is $T$-invariant, $z$ and $Ty$ must be orthogonal, hence $(z,Ty) = 0$. This means $(Tz,y) = 0$ for all $y \in Y$, hence $Tz \in Y^\perp$. Thus $Y^\perp$ is also $T$-invariant.
    \end{proof}

\section*{Week 8 Problems}
\begin{exercise}[manual-num={1}]
With the notation we had in class: $X:= \{f \in C^2([a,b]) \mid f \text{ satisfies BCs from RSL }\}$, and $Lu = -((pu')' + qu)$. Establish Lagrange's identity:
    \begin{equation*}
    \begin{split}
        uLv - vLu = -(p(uv' - vu'))'
    \end{split}
    \end{equation*}
\end{exercise}
    \begin{proof}
        Observe that:
            \begin{equation*}
            \begin{split}
                uLv - vLu 
                & = -u((pv')' + qv) + v((pu')' + qu)\\
                & = -u(p'v' + pv'' + qv) + v(p'u' + pu'' + qu) \\
                & = -up'v' - upv'' - uqv + vp'u' + vpu'' + vqu \\
                & = -up'v' - upv'' + vp'u' + vpu'' \\
                & = -up'v' - upv'' - pu'v' + pu'v' + vp'u' + vpu'' \\
                & = -(puv')' + (pvu')' \\
                & = -(puv' - pvu')' \\
                & = -(p(uv' - vu'))'. \qedhere
            \end{split}
            \end{equation*}
        
    \end{proof}

\newpage
\begin{exercise}[manual-num={4}]
Consider the RSL:
    \begin{equation*}
    \begin{split}
        X'' = -\lambda X, \h6 0<x<1, \h6X(0) = 0,\h6 X'(1) = aX(1),
    \end{split}
    \end{equation*}
where $a > 0$. Show, graphically, that is $a$ is large enough then there will be exactly one negative eigenvalue. What value must $a$ exceed to guarantee this?
\end{exercise}
    \begin{proof}
        If $\lambda < 0$, take $\lambda := -\mu^2$. Our RSL becomes:
            \begin{equation*}
            \begin{split}
                X'' - \mu X = 0.
            \end{split}
            \end{equation*}
        The solution to this differential equation is:
            \begin{equation*}
            \begin{split}
                X(x) = A \sinh(\mu x) + B \cosh(\mu x).
            \end{split}
            \end{equation*}
        Plugging in our first boundary condition yields:
            \begin{equation*}
            \begin{split}
                0 = X(0) = B.
            \end{split}
            \end{equation*}
        Now plugging in the second boundary condition gives:
            \begin{equation*}
            \begin{split}
                \mu A \cosh (\mu) = a A \sinh(\mu). 
            \end{split}
            \end{equation*}
        Simplifying and rearranging gives:
            \begin{equation*}
            \begin{split}
                \mu = a \tanh(\mu).
            \end{split}
            \end{equation*}
        From the following graph, if $a > 1$ then there will be exactly one negative eigenvalue.
        \begin{center}
        \begin{tikzpicture}[>=stealth]
            \begin{axis}[
                xmin=-7,xmax=7,
                ymin=-3,ymax=3,
                axis x line=middle,
                axis y line=middle,
                axis line style=<->,
                xlabel={$x$},
                ylabel={$y$},
                xtick=\empty,
                ytick=\empty,
                ]
                \addplot[no marks,black,-] expression[domain=-4:4,samples=100]{(1/2)*tanh(x)} 
                            node[pos=1,anchor=west]{\tiny$\frac{1}{2}\tanh(\mu)$}; 
                \addplot[no marks,black,-] expression[domain=-4:4,samples=100]{tanh(x)} 
                            node[pos=1,anchor=west]{\tiny$\tanh(\mu)$}; 
                \addplot[no marks,black,-] expression[domain=-4:4,samples=100]{(3/2)*tanh(x)} 
                            node[pos=1,anchor=west]{\tiny$\frac{3}{2}\tanh(\mu)$}; 
                \addplot[no marks,black,-] expression[domain=-4:4,samples=100]{2*tanh(x)} 
                            node[pos=1,anchor=west]{\tiny$2\tanh(\mu)$};             
                \addplot[no marks,black,-] expression[domain=-2.75:2.75,samples=100]{x} 
                            node[pos=1,anchor=west]{\tiny$\mu$};             
            \end{axis}
        \end{tikzpicture}
        \end{center}
    \end{proof}

\newpage
\begin{exercise}[manual-num={5}]
Consider the RSL:
    \begin{equation*}
    \begin{split}
        X'' = -\lambda X, \h6 0<x<1, \h6X'(0) = aX(0),\h6 X'(1) = -bX(1),
    \end{split}
    \end{equation*}
where $a,b \geq 0$. Without solving for the eigenvalues show that all eigenvalues must be greater than or equal to 0. (\textit{Hint: Let $L := -\frac{d^2}{dx^2}$; assuming $Lu = \lambda u$ for a nonzero function $u$ which satisfies the BCs, calculate $(Lu,u)$ in two different ways.})
\end{exercise}
    \begin{proof}
        Defining $L := -\frac{d^2}{dx^2}$, we have:
            \begin{equation*}
            \begin{split}
                (Lu,u)
                & = \int_0^1 -u''(x)\overline{u(x)}dx \\
                & = -u'(x)\overline{u(x)} \h3\bigg|_0^1 + \int_0^1 u'(x)\overline{u'(x)}dx \\
                & = -u'(1)\overline{u(1)}+u'(0)\overline{u(0)} + \int_0^1 u'(x)\overline{u'(x)}dx \\
                & = bu(1)\overline{u(1)} +au(0)\overline{u(0)} + \int_0^1 u'(x)\overline{u'(x)}dx \\ 
                & = b|u(1)|^2 + a|u(0)|^2 + \int_0^1 |u'(x)|^2 dx.
            \end{split}
            \end{equation*}
        Similarly:
            \begin{equation*}
            \begin{split}
                (\lambda u, u )
                & = \int_0^1 \lambda u(x) \overline{u(x)}dx \\
                & = \lambda \int_0^1 |u(x)|^2dx.
            \end{split}
            \end{equation*}
        Since $Lu = \lambda u$, we have $(Lu,u) = (\lambda u,u)$. Solving for $\lambda$ gives:
            \begin{equation*}
            \begin{split}
                \lambda = \frac{b|u(1)|^2 + a|u(0)|^2 + \int_0^1 |u'(x)|^2 dx}{\int_0^1 |u(x)|^2dx}.
            \end{split}
            \end{equation*}
        Since $u(x)$ is a nonzero function, the right hand side of the above equation is not undefined. Moreover, each term in the numerator is greater than or equal to zero, hence $\lambda \geq 0$.
    \end{proof}

\newpage
\begin{exercise}[manual-num={6}]
Let $k$ be a continuous function on $[a,b] \times [a,b]$. Define the operator $K:L^2([a,b]) \rightarrow L^2([a,b])$ by:
    \begin{equation*}
    \begin{split}
        (Kg)(x) = \int_a^b k(x,t)g(t)dt.
    \end{split}
    \end{equation*}
Use Cauchy-Schwartz and the continuity of $k$ to show that $Kg$ is continuous for any $g \in L^2([a,b])$.
\end{exercise}
    \begin{proof}
        Let $g \in L^2([a,b])$. Let $x \in [a,b]$ and $\epsilon > 0$. Since $k$ is continuous, there exists a $\delta > 0$ such that, for all $y \in [a,b]$, we have $|x-y| < \delta$ implies:
            \begin{equation*}
            \begin{split}
                |k(x,t) - k(y,t)| < \frac{\epsilon}{\lnorm g \rnorm_{L^2}}.
            \end{split}
            \end{equation*}
        For any $y \in [a,b]$, if $|x-y| < \delta$ then:
            \begin{equation*}
            \begin{split}
                |(Kg)(x) - (Kg)(y)| 
                & \leq \int_a^b |k(x,t) - k(y,t)|\cdot|g(t)|dt \\
                & \leq \left( \int_a^b |k(x,t) - k(y,t)|^2dt \right)^\frac{1}{2}\cdot \lnorm g \rnorm_{L^2} \\
                & < \left( \int_a^b \frac{\epsilon^2}{\lnorm g \rnorm_{L^2}^2} dt \right)^\frac{1}{2}\lnorm g \rnorm_{L^2} \\
                & = \frac{\epsilon}{\lnorm g \rnorm_{L^2}} \lnorm g \rnorm_{L^2} \\
                & = \epsilon.
            \end{split}
            \end{equation*}
        Thus $Kg$ is continuous.
    \end{proof}


\end{document}